\clearpage
\section{The philosophy of physical approaches to biological modeling}
\label{sec:philosophy}
% EDITORIAL ADDITION: navigation or conceptual bridge
The case studies reveal recurring structures: gain and loss, probability flow, state counting, and useful changes of coordinates.
The closing questions concern what these descriptions explain and where their assumptions stop being adequate.
% Source: day4.tex:16; unit d4.1
\subsection{The sameness of things}
\label{sec:sameness}


% ============================================================

The same physical ideas appear repeatedly across very different biological systems.
Random walks, for example, arise in processes ranging from molecular hopping to the conformations of polymers.
Likewise, similar ideas about binding, switching, cooperativity, and feedback appear in many important regulatory systems.

Examples discussed in the course include:

\begin{itemize}
    \item \textbf{Chemotaxis.}
    Measurements of bacterial chemotaxis reveal highly sensitive receptor responses, while models of interacting two-state receptors show how the same receptor clusters can operate in regimes dominated by sensitivity or cooperativity \cite{sourjik2002receptor,keymer2006chemosensing}.

    \item \textbf{Quorum sensing.}
    Measurements of the LuxN receptor in \textit{Vibrio harveyi} show how a quantitative receptor model can be used to infer signaling parameters from in vivo responses \cite{swem2008deducing}.

    \item \textbf{Gene expression.}
    Single-cell measurements of gene expression show how stochastic production and degradation of molecules generate fluctuations and transcriptional bursts \cite{golding2005real}.
\end{itemize}


% ============================================================

% EDITORIAL ADDITION: navigation or conceptual bridge
\paragraph{Would life look different elsewhere?}
\label{sec:elsewhere}
The distinction between general organizing principles and contingent molecular implementations was introduced in section~\ref{sec:life-features}.
After the case studies, the question can be revisited more concretely: which explanations depend on particular molecules, and which depend on structures such as autocatalysis, stochastic sampling, or energetic coupling?
% The opening comparison table is a thought experiment, not an established account of extraterrestrial life.
% Source: day4.tex:107; unit d4.2.2
\subsection{Defiance: maintaining living systems away from equilibrium}
\label{sec:defiance}

% ============================================================

One way to think about living systems is through their ability to maintain states that would not persist at thermodynamic equilibrium.
They do so by coupling internal processes to continual inputs of free energy, thereby breaking detailed balance (every microscopic transition is balanced by its reverse, so there are no net probability currents) and sustaining nonequilibrium dynamics \cite{gnesotto2018broken}.

Examples include:

\begin{itemize}
    \item creating and maintaining concentration gradients;
    \item increasing biochemical specificity through energy-driven kinetic proofreading \cite{hopfield1974kinetic};
    \item adapting to changes in the environment;
    \item enforcing ordered kinetic sequences, in which a later step cannot occur before a required earlier step (``socks before shoes'' kinetics);
    \item sustaining organized structures that persist only while energy is continually supplied and dissipated.
\end{itemize}

Living systems therefore do not violate thermodynamics.
They maintain nonequilibrium organization through continual free-energy input and dissipation \cite{gnesotto2018broken}.


% ============================================================

% Source: day5.tex:634; unit d5.6
\subsection{What physical approaches bring to biology}
\label{sec:ignorance}


% ============================================================

Physical approaches to biology are not defined only by the use of equations or quantitative measurements.
They also provide particular ways of deciding what constitutes an explanation or an understanding of a biological system, such as:

\begin{itemize}
    \item providing alternative definitions of what it means to understand a phenomenon;
    \item exposing hidden sameness between apparently unrelated biological systems;
    \item providing a different perspective on what is meant by a mechanism;
    \item allowing null hypotheses to be formulated explicitly and quantitatively;
    \item forcing us to state clearly what is known and what remains unknown, in the spirit of what Joshua Epstein calls ``militant ignorance'' \cite{epstein2008};
    \item emphasizing conceptual knowledge rather than the accumulation of isolated facts.
\end{itemize}

A recurring theme throughout the course is that the choice of representation, model, or null hypothesis is itself part of the scientific reasoning.
The aim is not merely to calculate a quantity, but to identify a description in which the essential structure of the biological problem becomes visible.

\subsubsection{Dimensionless numbers as quantitative perspective}

A final example of the physical approach is the construction of dimensionless numbers.
By comparing two effects or scales so that their units cancel, these ratios replace absolute magnitude with relative importance.
Values much smaller than one, of order one, or much larger than one identify regimes in which the denominator dominates, the two effects are comparable, or the numerator dominates, respectively.

A well-chosen dimensionless number can reveal which effects govern a system, which can be neglected, and which apparently different phenomena share the same underlying structure.
Its significance therefore lies less in its name or precise value than in the choice of quantities being compared.
In this way, dimensionless numbers exemplify a central contribution of physical thinking to biology: finding representations that make otherwise abstract relationships quantitatively intelligible.

The following dimensionless number illustrate this spirit. Some of these numbers aren't canonically named, but they are useful to think about. 

\paragraph{Biologically relevant examples}

\begin{center}
\renewcommand{\arraystretch}{1.8}
\begin{tabularx}{0.94\textwidth}{
    @{}
    >{\raggedright\arraybackslash}p{0.28\textwidth}
    |
    >{\raggedright\arraybackslash}X
    @{}
}
\hline
\textbf{Dimensionless number} & \textbf{Ratio} \\
\hline
Reynolds number &
$\mathrm{Re}
= \dfrac{\text{inertial force}}{\text{viscous force}}
= \dfrac{u\rho L}{\mu}$ 
\medskip
\\
Péclet number &
$\mathrm{Pe}
= \dfrac{\text{diffusion time}}{\text{advection time}}
= \dfrac{L^2/D}{L/V}$ 
\medskip
\\
Jesus number &
$\mathrm{Je}
= \dfrac{\text{surface-tension force}}{\text{gravitational force}}
= \dfrac{\gamma \ell}{\rho \ell^3 g}$ 
\medskip
\\
Langmuir number &
$\mathrm{La}
= \dfrac{\text{binding energy}}{\text{entropy loss of ligands}}
= \dfrac{\epsilon_b}{k_B T \mathrm{ln}\frac{c}{c_0}}$ 
\medskip
\\
Bjerrum number &
$\mathrm{Bj}
= \dfrac{\text{electrostatic energy}}{\text{thermal energy}}
= \dfrac{e^2}{4 \pi \epsilon_0 D \ell k_B T}$  
\medskip
\\
Hooke number &
$\mathrm{Ho}
= \dfrac{\text{elastic energy stored in a spring}}{\text{thermal energy}}
= \dfrac{k \ell^2}{2k_B T}$ 
\medskip
\\
\hline
\end{tabularx}
\end{center}

\paragraph{Anthropogenic examples}

\begin{center}
\renewcommand{\arraystretch}{1.8}
\begin{tabularx}{0.94\textwidth}{
    @{}
    >{\raggedright\arraybackslash}p{0.28\textwidth}
    |
    >{\raggedright\arraybackslash}X
    @{}
}
\hline
\textbf{Dimensionless number} & \textbf{Ratio} \\
\hline
% \medskip
Humanification number &
$\mathrm{Hu}
= \dfrac{\text{anthropomass}}{\text{biomass}}$
\medskip
\\
Keeling number &
$\mathrm{Ke}
= \dfrac{\text{CO$_2$ added by humans}}{\text{total CO$_2$ before humans}}$
\medskip
\\
Wildlife number &
$\mathrm{Wi}
= \dfrac{\text{biomass of humans + livestock}}{\text{biomass of wild animals}}$
\medskip
\\
Haber number &
$\mathrm{Ha}
= \dfrac{\text{N$_2$ flux through Haber-Bosch process}}{\text{N$_2$ flux through biological processes}}$
\medskip
\\
Amazon number &
$\mathrm{Am}
= \dfrac{\text{daily human water use}}{\text{Amazon drainage rate}}$
\medskip
\\
Mortality number &
$\mathrm{Mo}
= \dfrac{\text{accidental deaths in the US}}{\text{annual deaths in the US}}$
\medskip
\\
\hline
\end{tabularx}
\end{center}

% The significance of these numbers lies not in their names, nor in memorizing their precise values.
% It lies in the decision about what should be compared.
% Choosing a denominator turns an isolated measurement into a statement about relative importance, and a well-chosen dimensionless ratio can make an otherwise abstract scale immediately intelligible.
% This is one of the central contributions of physical thinking to biology:
% finding representations that reveal which effects dominate, which can be neglected, and which apparently different phenomena share the same underlying structure.

% \paragraph{Conclusion}
Across these examples, the physical approach does not provide a single theory of life, but a disciplined way of asking questions about living systems.
It begins by estimating scales, choosing variables and representations, stating assumptions, and constructing the simplest model capable of making a testable prediction.
Means, distributions, random walks, state counting, and changes of coordinates are useful not because living systems are simple, but because they make precise what would otherwise remain qualitative or ambiguous.
A model is therefore neither a miniature copy of biology nor an end in itself; it is a tool for exposing what is known, what has been assumed, and what remains unexplained.
This combination of quantitative clarity, conceptual economy, and explicit ignorance is the enduring contribution of physical thinking to biology.